\section{Limitaciones y amenazas a la validez}
\label{sec:disc_limitaciones}
\pendiente{Se analizan las amenazas a la validez interna, externa, de constructo y de conclusión estadística,
junto con las limitaciones del trabajo: escala de los modelos, costo, ausencia de reportes redactados por
especialistas, generalización fuera de distribución, anclajes y generadores no evaluados, y generación de
reportes sobre el test original.

Para cada resultado con el que se responden las preguntas y se evalúan los objetivos se indica si hubo
evidencia independiente de los datos con los que se tomaron las decisiones (HIBA, la partición disjunta por
lesión, semillas nuevas) o si se evaluó sobre el mismo material, en cuyo caso su respaldo queda como
exploratorio.

Recibe del Capítulo~\ref{cp:introduccion} las limitaciones de alcance que antes estaban en su sección de
alcance y delimitación (commit f520295): el trabajo se pronuncia sobre imágenes dermatoscópicas, y nada se
traslada a la histopatología ni a la fotografía clínica sin volver a medirlo; las métricas automáticas
establecen fidelidad diagnóstica, no validez clínica, que exigiría la lectura de especialistas; no se
desarrolla un producto de uso clínico ni se lo valida con pacientes; y la restricción de cómputo deja fuera
a los modelos de mayor tamaño, que podrían rendir mejor.}
